\documentclass[11pt,a4paper]{article}
\usepackage[T1,T2A]{fontenc}
\usepackage[a4paper, top=20mm, bottom=20mm, left=20mm, right=13mm]{geometry}
\usepackage[utf8]{inputenc}
\usepackage{cmap}
\usepackage[russian, english]{babel}
\usepackage{textcomp}
\usepackage{indentfirst}
\usepackage{rotating, graphicx}
\usepackage{mwe}
\usepackage{wrapfig}
\usepackage{hyperref}
\usepackage{caption}
\usepackage{multirow}
\usepackage{array}
\usepackage{makecell}
\usepackage{listings}
\usepackage{fancyvrb}
\usepackage{amsmath}
\usepackage{xcolor}

% ======================= listings: настройка =======================
\lstset{
    basicstyle=\ttfamily\footnotesize,
    breaklines=true,
    frame=single,
    numbers=left,
    numberstyle=\tiny,
    showstringspaces=false,
    captionpos=b,
    extendedchars=true,
    inputencoding=utf8,
    language=Python,
    keywordstyle=\color{blue},
    commentstyle=\color{green!60!black},
    stringstyle=\color{red!70!black},
    tabsize=4,
    columns=flexible,
    literate=
    {а}{{\cyra}}1 {б}{{\cyrb}}1 {в}{{\cyrv}}1 {г}{{\cyrg}}1
    {д}{{\cyrd}}1 {е}{{\cyre}}1 {ё}{{\cyryo}}1 {ж}{{\cyrzh}}1
    {з}{{\cyrz}}1 {и}{{\cyri}}1 {й}{{\cyrishrt}}1 {к}{{\cyrk}}1
    {л}{{\cyrl}}1 {м}{{\cyrm}}1 {н}{{\cyrn}}1 {о}{{\cyro}}1
    {п}{{\cyrp}}1 {р}{{\cyrr}}1 {с}{{\cyrs}}1 {т}{{\cyrt}}1
    {у}{{\cyru}}1 {ф}{{\cyrf}}1 {х}{{\cyrh}}1 {ц}{{\cyrc}}1
    {ч}{{\cyrch}}1 {ш}{{\cyrsh}}1 {щ}{{\cyrshch}}1 {ъ}{{\cyrhrdsn}}1
    {ы}{{\cyrery}}1 {ь}{{\cyrsftsn}}1 {э}{{\cyrerev}}1 {ю}{{\cyryu}}1
    {я}{{\cyrya}}1
    {А}{{\CYRA}}1 {Б}{{\CYRB}}1 {В}{{\CYRV}}1 {Г}{{\CYRG}}1
    {Д}{{\CYRD}}1 {Е}{{\CYRE}}1 {Ё}{{\CYRYO}}1 {Ж}{{\CYRZH}}1
    {З}{{\CYRZ}}1 {И}{{\CYRI}}1 {Й}{{\CYRISHRT}}1 {К}{{\CYRK}}1
    {Л}{{\CYRL}}1 {М}{{\CYRM}}1 {Н}{{\CYRN}}1 {О}{{\CYRO}}1
    {П}{{\CYRP}}1 {Р}{{\CYRR}}1 {С}{{\CYRS}}1 {Т}{{\CYRT}}1
    {У}{{\CYRU}}1 {Ф}{{\CYRF}}1 {Х}{{\CYRH}}1 {Ц}{{\CYRC}}1
    {Ч}{{\CYRCH}}1 {Ш}{{\CYRSH}}1 {Щ}{{\CYRSHCH}}1 {Ъ}{{\CYRHRDSN}}1
    {Ы}{{\CYRERY}}1 {Ь}{{\CYRSFTSN}}1 {Э}{{\CYREREV}}1 {Ю}{{\CYRYU}}1
    {Я}{{\CYRYA}}1
}

\newcounter{imgcnt}
\setcounter{imgcnt}{1}
\newcounter{codecnt}
\setcounter{codecnt}{1}

\begin{document}
\sloppy                            % <-- убирает Overfull \hbox
\emergencystretch=1em
\setlength{\footskip}{10mm}

% ==================== ТИТУЛЬНЫЙ ЛИСТ ====================
\newpage\thispagestyle{plain}
\thispagestyle{empty}
\begin{center}
\textbf{МИНИСТЕРСТВО НАУКИ И ВЫСШЕГО ОБРАЗОВАНИЯ\\
РОССИЙСКОЙ ФЕДЕРАЦИИ\\[2ex]
\mbox{ФЕДЕРАЛЬНОЕ ГОСУДАРСТВЕННОЕ БЮДЖЕТНОЕ ОБРАЗОВАТЕЛЬНОЕ}\\
УЧРЕЖДЕНИЕ ВЫСШЕГО ОБРАЗОВАНИЯ\\
«ВЯТСКИЙ ГОСУДАРСТВЕННЫЙ УНИВЕРСИТЕТ»\\[2ex]
Институт математики и информационных систем\\[2ex]
Факультет автоматики и вычислительной техники\\[2ex]
Кафедра электронных вычислительных машин}
\end{center}
\vspace{35mm}

\begin{center}
Отчёт по проекту\\
<<Weather Forecast>>\\
по дисциплине\\
<<Программирование>>\\[1ex]
\end{center}
\vspace{45mm}
\begin{flushleft}
Выполнил студент гр. ИВТб-3302-05-00 \hspace{42mm} \rule[-0,5mm]{30mm}{0.15mm}\,/Иванов И.И./\\
\mbox{\tiny\hspace{124mm}(Подпись)}\\[1ex]
Проверил профессор кафедры ЭВМ \hspace{45mm} \rule[-0,5mm]{30mm}{0.15mm}\,/Петров П.П./\\
\mbox{\tiny\hspace{122mm}(Подпись)}
\end{flushleft}
\vfill
\begin{center}
Киров\\
2026
\end{center}

% ==================== ОСНОВНАЯ ЧАСТЬ ====================
\newpage
\topmargin = -35pt

\section*{\normalsize Цель работы}

Разработать приложение для автоматического получения прогноза погоды, сохранения его в базе данных и формирования Markdown-отчёта.

\section*{\normalsize Задание}

Изучить и описать архитектуру приложения \texttt{Weather Forecast}, реализованного на языке Python. Приложение должно:
\begin{itemize}
    \item определять текущую локацию пользователя по IP-адресу;
    \item запрашивать прогноз погоды на 4 дня через OpenWeatherMap API;
    \item агрегировать 3-часовые данные в дневные записи;
    \item сохранять данные в SQLite без дубликатов;
    \item формировать отчёт в формате Markdown;
    \item выводить сводку в консоль.
\end{itemize}

\section*{\normalsize Теоретические сведения}

В проекте используются следующие технологии и библиотеки:
\begin{itemize}
    \item \textbf{Python 3.10+} --- язык программирования, поддерживающий современные конструкции (аннотации типов, \texttt{dataclasses}).
    \item \textbf{SQLAlchemy 2.0} --- ORM для работы с базой данных SQLite. Позволяет описывать модели как классы Python и выполнять запросы через сессии.
    \item \textbf{requests} --- библиотека для выполнения HTTP-запросов к внешним API.
    \item \textbf{python-dotenv} --- загрузка переменных окружения из файла \texttt{.env}.
    \item \textbf{OpenWeatherMap API} --- внешний сервис для получения прогноза погоды. Используется endpoint \texttt{/data/2.5/forecast}, возвращающий данные с шагом 3 часа.
    \item \textbf{IP-геолокация} --- сервис \texttt{ip-api.com} для определения города и координат по IP-адресу.
    \item \textbf{Markdown} --- текстовый формат для формирования отчёта.
\end{itemize}

\section*{\normalsize Ход работы}

\subsection*{Архитектура приложения}

Приложение состоит из нескольких модулей, каждый из которых отвечает за свою часть функциональности:
\begin{itemize}
    \item \texttt{env\_manager.py} --- загрузка и валидация переменных окружения, настройка логирования.
    \item \texttt{geoloc.py} --- определение локации (IP-геолокация или фолбэк по городу).
    \item \texttt{weather\_api.py} --- запрос прогноза и агрегация данных по дням.
    \item \texttt{db\_manager.py} --- подключение к БД, модель \texttt{Weather}, сохранение записей.
    \item \texttt{report.py} --- формирование Markdown-отчёта.
    \item \texttt{main.py} --- точка входа, управляющая последовательностью вызовов.
\end{itemize}

Взаимодействие модулей можно описать следующей схемой:
\begin{center}
\texttt{main} $\rightarrow$ \texttt{geoloc} $\rightarrow$ \texttt{weather\_api} $\rightarrow$ \texttt{db\_manager} $\rightarrow$ \texttt{report}
\end{center}

\subsection*{Модуль env\_manager.py}

Модуль загружает переменные из \texttt{.env} с помощью \texttt{python-dotenv} и предоставляет их остальным частям программы. Обязательные переменные проверяются функцией \texttt{\_require}. Также настраивается базовое логирование.

\begin{lstlisting}[caption={Фрагмент env\_manager.py}, label={lst:env}]
import logging
import os
from dotenv import load_dotenv

load_dotenv()

LEVEL = os.getenv("LOG_LEVEL", "WARNING").upper()
logging.basicConfig(level=LEVEL,
                    format="%(asctime)s [%(levelname)s] %(name)s: %(message)s")
log = logging.getLogger(__name__)

def _require(name: str) -> str:
    v = os.getenv(name)
    if v is None:
        raise RuntimeError(f"{name} is not set")
    return v

DB_URL: str = _require("DB_URL")
DB_USER: str = os.getenv("DB_USER", "")
DB_PASSWORD: str = os.getenv("DB_PASSWORD", "")
WEATHER_API_KEY: str = _require("WEATHER_API_KEY")
WEATHER_API_URL: str = _require("WEATHER_API_BASE_URL")
GEO_API_URL: str | None = os.getenv("GEO_API_URL")
GEO_API_FALLBACK_CITY: str | None = os.getenv("GEO_API_FALLBACK_CITY")
OUTPUT_FILE: str = os.getenv("OUTPUT_FILE", "weather_report.md")
\end{lstlisting}

\subsection*{Модуль geoloc.py}

Определяет локацию пользователя. Если задан \texttt{GEO\_API\_URL}, выполняется запрос к сервису IP-геолокации. При успехе возвращается объект \texttt{Location} с городом и координатами. В случае ошибки или отсутствия URL используется запасной город из \texttt{GEO\_API\_FALLBACK\_CITY}.

\begin{lstlisting}[caption={Функция get\_location в geoloc.py}, label={lst:geo}]
@dataclass
class Location:
    city: str
    lat: float | None = None
    lon: float | None = None

def get_location() -> Location:
    if GEO_API_URL:
        try:
            r = requests.get(GEO_API_URL, timeout=5)
            r.raise_for_status()
            data = r.json()
            city = data.get("city")
            lat = data.get("lat") or data.get("latitude")
            lon = data.get("lon") or data.get("longitude")
            if city and lat is not None and lon is not None:
                return Location(city, float(lat), float(lon))
        except requests.RequestException as e:
            log.warning("Request to %s failed: %s", GEO_API_URL, e)
        except ValueError as e:
            log.warning("JSON parse error: %s", e)
    if GEO_API_FALLBACK_CITY:
        return Location(GEO_API_FALLBACK_CITY)
    raise RuntimeError("Could not determine city.")
\end{lstlisting}

\subsection*{Модуль weather\_api.py}

Содержит две основные функции: \texttt{fetch\_forecast} для запроса к OpenWeatherMap и \texttt{aggregate\_daily} для агрегации 3-часовых данных в дневные.

\begin{lstlisting}[caption={Функции fetch\_forecast и aggregate\_daily}, label={lst:weather}]
def fetch_forecast(loc: Location) -> list[dict]:
    params = {"appid": WEATHER_API_KEY, "units": "metric", "lang": "ru"}
    if loc.lat is not None and loc.lon is not None:
        params["lat"] = loc.lat
        params["lon"] = loc.lon
    else:
        params["q"] = loc.city
    r = requests.get(f"{WEATHER_API_URL}/forecast", params=params, timeout=10)
    if r.status_code == 404:
        raise RuntimeError(f"City not found: {loc.city}")
    r.raise_for_status()
    return r.json()["list"]

def aggregate_daily(items: list[dict], days: int = 4) -> list[dict]:
    by_day = defaultdict(list)
    for item in items:
        d = datetime.fromisoformat(item["dt_txt"]).date()
        by_day[d].append(item)
    today = datetime.now(tz=timezone.utc).date()
    result = []
    for d in sorted(by_day):
        if d < today:
            continue
        if len(result) >= days:
            break
        entries = by_day[d]
        temps_min = [e["main"]["temp_min"] for e in entries]
        temps_max = [e["main"]["temp_max"] for e in entries]
        humidity = [e["main"]["humidity"] for e in entries]
        wind = [e["wind"]["speed"] for e in entries]
        descriptions = [e["weather"][0]["description"] for e in entries]
        result.append({
            "date": d,
            "temp_min": min(temps_min),
            "temp_max": max(temps_max),
            "humidity": int(sum(humidity) / len(humidity)),
            "wind_speed": max(wind),
            "description": max(set(descriptions), key=descriptions.count),
        })
    return result
\end{lstlisting}

\subsection*{Модуль db\_manager.py}

Отвечает за подключение к базе данных, описание модели \texttt{Weather} и сохранение записей. Используется уникальное ограничение по паре \texttt{(city, forecast\_date)}, чтобы избежать дублирования.

\begin{lstlisting}[caption={Модель Weather и функция save\_forecast}, label={lst:db}]
class Weather(Base):
    __tablename__ = "weather_forecast"
    __table_args__ = (
        UniqueConstraint("city", "forecast_date", name="uq_city_forecast_date"),
    )
    id = Column(Integer, primary_key=True)
    city = Column(String(100), nullable=False)
    forecast_date = Column(Date, nullable=False)
    temp_min = Column(Float)
    temp_max = Column(Float)
    humidity = Column(Integer)
    wind_speed = Column(Float)
    description = Column(String(100))
    created_at = Column(DateTime)

def save_forecast(city: str, rows: list[dict]) -> list[Weather]:
    existing = {d for (d,) in session.query(Weather.forecast_date)
                .filter(Weather.city == city).all()}
    saved = []
    for row in rows:
        if row["date"] in existing:
            continue
        w = Weather(city=city, forecast_date=row["date"],
                    temp_min=row["temp_min"], temp_max=row["temp_max"],
                    humidity=row["humidity"], wind_speed=row["wind_speed"],
                    description=row["description"],
                    created_at=datetime.now(tz=timezone.utc))
        session.add(w)
        saved.append(w)
    session.commit()
    return saved
\end{lstlisting}

\subsection*{Модуль report.py}

Формирует Markdown-отчёт с выравниванием колонок. Функция \texttt{\_render} создаёт таблицу, а \texttt{write\_markdown} записывает её в файл.

\begin{lstlisting}[caption={Функция \_render в report.py}, label={lst:report}]
HEADERS = ["Date", "Min T (C)", "Max T (C)", "Description", "Humidity (%)", "Wind (m/s)"]

def _row_from(r) -> list[str]:
    return [str(r.forecast_date), str(round(r.temp_min)), str(round(r.temp_max)),
            r.description.capitalize(), str(r.humidity), str(round(r.wind_speed))]

def _render(city: str, rows: list) -> str:
    if not rows:
        return f"# Weather forecast\n\nLocation: {city}\n\nNo data.\n"
    body = [_row_from(r) for r in rows]
    widths = [len(h) for h in HEADERS]
    for row in body:
        for i, cell in enumerate(row):
            widths[i] = max(widths[i], len(cell))
    def line(cells):
        return "| " + " | ".join(c.ljust(widths[i])
                                 for i, c in enumerate(cells)) + " |"
    separator = "|" + "|".join("-" * (w + 2) for w in widths) + "|"
    period = f"{rows[0].forecast_date} - {rows[-1].forecast_date}"
    out = ["# Weather forecast", f"Location: {city}",
           f"Period: {period}", "", line(HEADERS), separator]
    out.extend(line(row) for row in body)
    return "\n".join(out) + "\n"
\end{lstlisting}

\subsection*{Главный модуль main.py}

Управляет последовательностью действий: определение локации, запрос прогноза, агрегация, сохранение, выборка из БД, формирование отчёта и вывод сводки.

\begin{lstlisting}[caption={Функция main в main.py}, label={lst:main}]
def main() -> int:
    loc = get_location()
    raw = fetch_forecast(loc)
    daily = aggregate_daily(raw, days=4)
    saved = save_forecast(loc.city, daily)
    today = datetime.now(tz=timezone.utc).date()
    period_end = today + timedelta(days=3)
    all_rows = session.query(Weather).filter(
        Weather.city == loc.city,
        Weather.forecast_date >= today,
        Weather.forecast_date <= period_end
    ).order_by(Weather.forecast_date).all()
    write_markdown(loc.city, all_rows)
    period_start = all_rows[0].forecast_date if all_rows else "-"
    actual_period_end = all_rows[-1].forecast_date if all_rows else "-"
    print_summary(loc, period_start, actual_period_end, len(saved), OUTPUT_FILE)
    return 0
\end{lstlisting}

\subsection*{Настройка и запуск}

Для работы приложения необходимо создать файл \texttt{.env} на основе \texttt{.env.example} и заполнить его. Основные переменные приведены в таблице~\ref{tab:env}.

\begin{table}[!ht]
\centering
\captionsetup{justification=raggedleft,singlelinecheck=false}
\caption*{Таблица 1}
\captionsetup{justification=centering,singlelinecheck=false}
\caption*{Переменные окружения}
\label{tab:env}
\renewcommand{\arraystretch}{1.12}
\begin{tabular}{|>{\ttfamily}p{0.35\textwidth}|>{\raggedright\arraybackslash}p{0.55\textwidth}|}
\hline
\textbf{Переменная} & \textbf{Описание}\\
\hline
DB\_URL & URL базы данных (например, sqlite:///./weather.db)\\
\hline
DB\_USER & Имя пользователя БД (для SQLite пусто)\\
\hline
DB\_PASSWORD & Пароль БД (для SQLite пусто)\\
\hline
WEATHER\_API\_KEY & API-ключ OpenWeatherMap (обязателен)\\
\hline
WEATHER\_API\_BASE\_URL & Базовый URL API (https://api.openweathermap.org/data/2.5)\\
\hline
GEO\_API\_URL & URL сервиса IP-геолокации (https://ip-api.com/json/)\\
\hline
GEO\_API\_FALLBACK\_CITY & Запасной город, если IP-геолокация не сработала\\
\hline
LOG\_LEVEL & Уровень логирования (DEBUG, INFO, WARNING, ERROR, CRITICAL)\\
\hline
OUTPUT\_FILE & Путь к Markdown-файлу отчёта\\
\hline
\end{tabular}
\end{table}

Зависимости проекта перечислены в \texttt{requirements.txt}:
\begin{itemize}
    \item requests 2.34.2
    \item SQLAlchemy 2.0.54
    \item python-dotenv 1.2.3
\end{itemize}

Запуск производится командой \texttt{python src/main.py} или через скрипт \texttt{run.sh}.

\subsection*{Пример вывода}

Пример Markdown-отчёта:
\begin{verbatim}
# Weather forecast
Location: Kirov
Period: 2026-09-28 - 2026-10-01

| Date       | Min T (C) | Max T (C) | Description | Humidity (%) | Wind (m/s) |
|------------|-----------|-----------|-------------|--------------|------------|
| 2026-09-28 | 8         | 9         | Clear       | 91           | 2          |
| 2026-09-29 | 5         | 16        | Clear       | 77           | 3          |
| 2026-09-30 | 4         | 16        | Clear       | 78           | 2          |
| 2026-10-01 | 6         | 15        | Cloudy      | 77           | 3          |
\end{verbatim}

Пример консольной сводки:
\begin{verbatim}
==================================================
  City:              Frankfurt am Main
  Coordinates:       50.0827, 8.6297
  Forecast period:   2026-09-28 - 2026-10-01
  Saved records:     0
  Report file:       weather_report.md
==================================================
\end{verbatim}

\section*{\normalsize Вывод}

В ходе выполнения проекта разработано приложение \texttt{Weather Forecast}, демонстрирующее модульный подход к построению программ на Python. Освоены работа с внешними API (OpenWeatherMap, ip-api.com), ORM SQLAlchemy для взаимодействия с SQLite, обработка ошибок и логирование, а также формирование отчётов в формате Markdown. Приложение успешно решает задачу автоматического получения и сохранения прогноза погоды.

\end{document}
